\section{Memory Models}

\begin{definition}
    \textbf{Memory Reordering:} Compilers/hardware may reorder memory operations as long as \textit{sequential} semantics hold (optimizes registries, caching, dead code). May break concurrency (e.g., hoisted checks $\implies$ infinite loops). L1 caches are private, writes may not immediately publish.
\end{definition}

\begin{remark}
    \textbf{\texttt{volatile}:} Prevents local caching/reordering; ensures immediate visibility across threads. Memory barriers strictly enforce ordering.
\end{remark}

\begin{important}
    Volatile arrays only guarantee that the array reference is volatile, not the elements.
\end{important}

\begin{remark}
    Almost all common architectures allow the reordering of \textbf{Stores after Loads} to maximize performance, unless explicitly prevented by memory barriers.
\end{remark}

\subsection{Java Memory Model (JMM)}
Contract providing minimal memory operation guarantees.

\begin{definition} 
    \textbf{Program Order (PO):} Total intra-thread execution trace order (no global sync).
\end{definition}

\begin{definition}
    \textbf{Synchronization Actions (SA):} E.g., Read/Write \texttt{volatile}, lock acquire/release.
\end{definition}

\begin{remark}
    SA establish HB dependencies between threads:
    \begin{itemize}
        \item Write to volatile in thread $A$ $\rightarrow$ Read from volatile in thread $B$
        \item Lock release in thread $A$ $\rightarrow$ Lock acquire in thread $B$
    \end{itemize}
\end{remark}

\begin{definition}
    \textbf{Synchronization Order (SO):} Total global order of all SA. All threads see these actions in the same global order.
\end{definition}

\begin{definition}
    \textbf{Happens-Before (HB):} Transitive closure of PO \& Synchronizes-With (SW). $A \xrightarrow{HB} B \implies$ $B$ sees results of $A$.
\end{definition}

\begin{remark}
    \textbf{HB-Consistency \& Data Races:} JMM permits data races (unsynchronized concurrent access where $\ge 1$ is a write). Unsynchronized reads may observe stale values or non-HB ordered writes.
\end{remark}
